\section{varFDTD Verification}
varFDTD is a reduced 2.5D time-domain method that uses an effective-index approximation to collapse the 3D SOI geometry into a 2D propagation problem. It is less expensive than full 3D FDTD, while still providing broadband time-domain verification of the final microring layout \cite{labmanual}.

The intended verification outputs are the final device geometry, field profile, through/drop monitor spectra, and effective-index monitor data. The file \path{models_original/ring_resonator_studentv2_varfdtd.lms} is available locally, but no exported varFDTD field image, spectrum, or monitor data was found. Therefore, the comparison table below is structured for completion after simulation export.

\begin{table}[H]
\centering
\caption{Comparison of abstraction levels. Only target-derived expressions are filled; missing simulation outputs are marked explicitly.}
\label{tab:method-comparison}
\begin{tabular}{lllll}
\toprule
Method & Key input & FSR (nm) & \(Q\) & Finesse \\
\midrule
INTERCONNECT & Circuit parameters & TODO & TODO & TODO \\
FDE theory & \(n_g\) and radius & 26.2 target & N/A & N/A \\
varFDTD & Final device geometry & TODO & TODO & TODO \\
\bottomrule
\end{tabular}
\end{table}

After varFDTD completion, the most important checks are whether the drop-port spectrum reproduces the target FSR near \SI{26.2}{\nano\metre}, whether the resonance linewidth gives the intended \(Q\), and whether the field profile confirms that power is coupled into the ring and transferred to the drop port at resonance.
